\subsection{NOTATION}

Given a real function \(f \in \mathcal{H}(\mathfrak{F}, \mathbf{R})\) it is often convenient in a formula to write \(O(f)\) for a function dominated by \(f\) , and \(o(f)\) for a negligible function relative to \(f\) . When, in a proof, there feature several functions dominated by the same function \(f\) (resp. negligible relative to \(f\) ) we denote them by \(O_1(f)\) , \(O_2(f)\) , etc. (resp. \(o_1(f)\) , \(o_2(f)\) , etc.).

Many authors write \(O(f)\) (resp. \(o(f)\) ) indiscriminately for all functions in a proof that are dominated by f (resp. negligible relative to f), an abuse of language which may risk confusion.

With this notation we can express props. 1, 2, 3 (V, p.~213) as follows: if \(g = O_1(f)\) and \(h = O(g)\) then \(h = O_2(f)\) ; one can write

\[
\sum_ {i = 1} ^ {n} \lambda_ {i} O _ {i} (f) = O _ {n + 1} (f) \quad (\lambda_ {i} \text {   scalars })\tag{1}
\]

\[
O (f) O (g) = O (f g).\tag{2}
\]

Similarly, prop. 4 (V, p.~215) shows that if \(g = O_1(f)\) and \(h = o(g)\) (resp. \(g = o_1(f)\) and \(h = O(g)\) then \(h = o_2(f)\) , and props. 5 and 6 (V, p.~5) can be expressed in the form

\[
\sum_ {i = 1} ^ {n} \lambda_ {i} o _ {i} (f) = o _ {n + 1} (f) \quad (\lambda_ {i} \text {   scalars })\tag{3}
\]

\[
o (f) O (g) = o (f g).\tag{4}
\]

The relation \(f \sim g\) is equivalent to \(f = g + o(g)\) . The notation \(O(1)\) (resp. \(o(1)\) ) denotes a function bounded on a set in \(\mathfrak{F}\) (resp. a function tending to 0 along \(\mathfrak{F}\) ).

